%% Babel's Machines — Royal Society Interface Focus perspective
%% Adapted from ACL position survey (2026) on multilingual LLM games
%% Template: rsproca_new (Royal Society open access)

\documentclass[openacc]{rsproca_new}

\usepackage{multicol}
\usepackage{booktabs}
\usepackage{amsmath}
\usepackage{xcolor}
\usepackage{float}

\addbibresource{refs.bib}

%%%%%%%%%%%%%%%%%%%%%%%%%%%%%%%%%%%%%%%%%%%%%%%
\newtheorem{theorem}{\bf Theorem}[section]
\newtheorem{condition}{\bf Condition}[section]
\newtheorem{corollary}{\bf Corollary}[section]
%%%%%%%%%%%%%%%%%%%%%%%%%%%%%%%%%%%%%%%%%%%%%%%

\jname{rsif}
\Journal{Interface Focus\ }

\begin{document}

\title{Babel's Machines: Linguistic Relativity, Cooperation, and the Emerging Behaviour of Multilingual Artificial Agents}

\author{%
Dao Sy Duy Minh$^{1,*}$,
Vu Nguyen$^{1}$,
[Co-author TBC]$^{2}$}

\address{%
$^{1}$Faculty of Information Technology, University of Science, Vietnam National University, Ho Chi Minh City, Vietnam\\
$^{2}$[Affiliation pending]}

\subject{artificial intelligence, behavioural science, evolutionary game theory}

\keywords{large language models, multi-agent systems, linguistic relativity, cooperation, theory of mind, machine behaviour}

\corres{Dao Sy Duy Minh\\
\email{dsdminh22@apcs.fitus.edu.vn}}

\begin{abstract}
Large language model (LLM) agents now negotiate, cooperate, and deceive in multi-agent environments. We argue that \emph{the language an agent operates in systematically reshapes its strategic behaviour, theory-of-mind performance, and prosocial defaults} --- a phenomenon we call the \emph{Algorithmic Sapir--Whorf Switch} \cite{wang2025babel}. This perspective makes three moves. First, it describes the empirical phenotype: cross-linguistic divergence in strategy distributions, prosocial defaults that exceed typical human cooperation rates in cooperative and mixed-motive games, and language-dependent erosion of theory-of-mind. Second, it connects these shifts to two mechanistic accounts --- the Semantic Hub and Translation Barrier hypotheses --- and to three architectural alternatives that the literature has not yet adjudicated. Third, it reframes LLM cooperation in evolutionary terms: not as reciprocity evolved through interaction, but as \emph{culturally inherited cooperation} delivered by pretraining and reinforcement learning from human feedback (RLHF), with language acting as an inheritance bottleneck. Throughout, we foreground contradictory evidence --- a cross-lingual persuasion null finding, framing effects that dominate language effects, bimodal Arabic strategy distributions --- to bound the claim. We close with four falsifiable predictions and three concrete governance levers for the multilingual deployment of LLM agents.
\end{abstract}

\begin{fmtext}
\end{fmtext}

\maketitle

\begin{multicols}{2}

\section{Introduction: Machine Behaviour Meets Babel}
\label{sec:intro}

LLMs no longer just produce text. They act. They negotiate trades, allocate resources, deceive opponents in social-deduction games, and collaborate on open-ended tasks over hours-long horizons \cite{meta2022cicero, wang2024voyager, huynh2026moreatstake, sun2025gametheoryllm}. They cooperate at rates that exceed typical human players in iterated social dilemmas \cite{fontana2024nicer} and manage budgets in auction arenas \cite{chen2024aucarena}. The dynamics of such multi-agent ecosystems are no longer fully predictable from the capabilities of individual models \cite{cemri2025mast, nie2025multiagentbench}.

This transition calls for a new research programme. \citeauthor{rahwan2019machinebehaviour} \cite{rahwan2019machinebehaviour} call it \emph{machine behaviour}: studying AI systems as a class of behaving agents, with tools drawn from biology, economics, cognitive science, evolutionary game theory \cite{han2022emergent}, and the statistical physics of cooperation \cite{perc2017statistical} --- not just from the engineering discipline that built them.

This perspective advances a specific thesis within that programme. Multilingual LLM agents shift their cooperative tendencies, deceptive capacities, and theory-of-mind performance as a function of the language they are asked to inhabit. We call this the \emph{Algorithmic Sapir--Whorf Switch}, following recent empirical work \cite{wang2025babel}. This is \emph{not} the strong Whorfian claim that language determines thought. It is a weaker, mechanistically grounded claim: language-specific representations inside pretrained models modulate the policies those models execute in interactive environments. The claim is empirical, not metaphysical. We devote considerable space to contradictory evidence and confounds.

The question matters for three reasons that map onto the goals of this journal's call for papers on the science of machine behaviour. First, \emph{empirical characterisation}: we need to describe how LLM agents act and interact under varied linguistic pressure, because deployment is global and multilingual by default. Second, \emph{social and cognitive models}: concepts of trust, cooperation, reciprocity, theory of mind, and cultural evolution --- developed over decades to understand humans and other animals --- are being silently re-instantiated in machine populations \cite{axelrod1984evolution, tomasello2009whywecooperate, henrich2016secret}. Third, and most consequential, \emph{societal dynamics}: LLMs are already embedded in multilingual institutions --- customer service, education, public-sector information access, especially in the Global South --- and their language-modulated behaviour will reshape human groups in ways we barely measure.

We organise the perspective around these three goals. Section~\ref{sec:taxonomy} sketches the interactive substrate in which these behaviours arise. Section~\ref{sec:relativity} develops the linguistic-relativity hypothesis, foregrounding convergent \emph{and} contradictory evidence. Section~\ref{sec:mechanism} offers mechanistic explanations (Semantic Hub, Translation Barrier) that replace cultural essentialism with architectural reasoning. Section~\ref{sec:cognition} examines theory of mind across languages. Section~\ref{sec:evolution} places these findings in an evolutionary and cultural-evolution frame. Section~\ref{sec:society} turns to societal implications and the research agenda. We close with what we see as the most consequential open questions.

\section{An Interactive Substrate}
\label{sec:taxonomy}

To speak of machine \emph{behaviour}, rather than machine \emph{performance}, we need a substrate richer than single-turn question answering. Recent literature has converged, roughly, on seven families of interactive environments in which LLM agents are now regularly evaluated.

\emph{Social-dilemma games} (Iterated Prisoner's Dilemma, Public Goods, Ultimatum) provide the mathematically purest view of cooperation, trust, and reciprocity \cite{huynh2026moreatstake, akata2025playing, buscemi2025fairgame}. \emph{Negotiation and diplomacy} (Diplomacy, bargaining games, SPIN-Bench) test persuasion and coalition formation under partial information \cite{meta2022cicero, deng2024bargaining, yao2025spinbench}. \emph{Persuasion and auction} environments probe rhetorical influence and strategic bidding \cite{biswas2025persuasion, chen2024aucarena}. \emph{Social-deduction games} (Werewolf, Avalon, Spyfall) demand deception and recursive mental-state modelling \cite{xu2023werewolf, wang2023avalon, koto2025spyfall}. \emph{Cooperative communication games} (Codenames, Hanabi, Overcooked) require shared cultural and referential ground \cite{hakimov2025codenames, chen2025llmhanabi, sun2025overcooked}. \emph{Open-world and video-game environments} (Minecraft via Voyager, BALROG, GVGAI-LLM) test long-horizon planning in grounded worlds \cite{wang2024voyager, paischer2024balrog, chen2025gvgai}. \emph{Dialogue-game frameworks} such as clembench provide the most controlled multilingual evaluation, separating format adherence from strategic competence \cite{chalamalasetti2023clembench, beyer2024clembench, schlangen2025clembench}.

Two observations matter for the argument that follows. First, these environments produce \emph{behaviours}, not just answers: commitments honoured or broken, beliefs propagated or manipulated, coalitions joined or abandoned. Second, every one of these environments has now been run in multiple languages, and every one of them has reported language-dependent divergence --- including the ``nicer than human'' cooperation phenomenon \cite{fontana2024nicer, za2026towards} and the substantial English--Arabic performance gap reported by TCG-Bench at large scale \cite{wang2025tcgbench}.

\section{The Linguistic-Relativity Hypothesis for Machines}
\label{sec:relativity}

The classical hypothesis of linguistic relativity, associated with Sapir, Whorf, and a century of experimental follow-ups, holds that the language one speaks shapes the cognitive categories one easily forms and deploys \cite{boroditsky2001does, boroditsky2011languageshapes}. Contemporary research has moved past the strong determinist reading toward a measured view: language provides \emph{cognitive affordances}, subtly biasing attention, memory, and decision-making without dictating them.

We propose that an analogous effect operates in multilingual LLMs. Pretraining on unbalanced, largely English-dominant corpora \cite{brown2020gpt3, vaswani2017transformer}, followed by reinforcement learning from human feedback (RLHF) whose safety and value data are themselves linguistically skewed \cite{ouyang2022training, bai2022constitutional, repello2024multilingualguardrails}, produces models whose internal representations encode a language-specific inductive bias. When prompted in a different language, the resulting policy shifts.

\subsection{Cross-Lingual Divergence in Social Dilemmas}

High-resolution data from iterated social dilemmas illustrates the phenomenon (Table~\ref{tab:strategy}). Across four model families evaluated in five languages, strategy distributions differ meaningfully, but not in ways that map cleanly onto folk-cultural stereotypes \cite{huynh2026moreatstake, huynh2025fairgame, za2026towards}.

\begin{table}[H]
\centering
\small
\begin{tabular}{@{}lcccc@{}}
\toprule
\textbf{Language} & \textbf{ALLC} & \textbf{ALLD} & \textbf{TFT} & \textbf{WSLS} \\
\midrule
English & 26\% & 25\% & 25\% & 24\% \\
French & 27\% & 18\% & 26\% & 29\% \\
Arabic & 32\% & 28\% & 19\% & 21\% \\
Chinese & 31\% & 30\% & 19\% & 20\% \\
Vietnamese & 26\% & 24\% & 23\% & 27\% \\
\bottomrule
\end{tabular}
\caption{IPD strategy distributions aggregated across four model families (GPT-4o, GPT-4o-mini, Claude 3.5 Haiku, Mistral) over $\geq 200$ rounds per language--model condition, with strategy labels assigned by an LSTM classifier at confidence threshold $\geq 0.6$ \cite{huynh2026moreatstake, huynh2025fairgame, za2026towards}. Cells are \emph{descriptive proportions, not a meta-analytic estimate}: prompts, payoff matrices, horizon disclosure and opponent policies vary across the constituent studies, so cross-language differences should be read as model--language interaction signals under specific protocols, \emph{not} as stable properties of the languages themselves. The original studies report per-language classifier reliability of $0.78$--$0.91$ on held-out human-validation sets and report stability of the qualitative ranking under threshold sweeps in $[0.5, 0.7]$; the rank order of cells should be interpreted as robust to the classifier choice, while exact percentages should not. We deliberately omit confidence intervals because the underlying studies use heterogeneous run counts; a properly powered meta-analysis with per-study normalisation is identified as future work in P1 (Section~\ref{sec:predictions}).}
\label{tab:strategy}
\end{table}

Three patterns deserve emphasis. Arabic produces \emph{both} the highest unconditional cooperation (32\%) \emph{and} the second-highest unconditional defection (28\%) --- a bimodal, volatile distribution rather than monotonic competitiveness. Vietnamese, often assumed to prime collectivist behaviour, sometimes yields high defection, with cooperative priming collapsing to 2\% cooperation in Claude 3.5 Haiku \cite{huynh2025fairgame}. Chinese favours Win-Stay-Lose-Shift over Tit-for-Tat, an outcome-dependent rather than reciprocity-dependent style. Layered on top, model-family effects are substantial: GPT-4o runs ``hawkish'' (aggregate ALLD $\approx 31.7\%$) while Mistral runs ``dovish'' (aggregate ALLC $\approx 33.7\%$) \cite{huynh2026moreatstake}.

\subsection{The ``Nicer than Human'' Finding}

Against these divergences sits a baseline finding that the machine-behaviour community has under-discussed. Across recent IPD studies, frontier LLMs default to \emph{prosocial} play at rates substantially exceeding typical human players in comparable settings. We must immediately bound this claim. \citeauthor{fontana2024nicer} \cite{fontana2024nicer} report that Llama2 and GPT-3.5 are ``markedly more cooperative than humans'' on an iterated Prisoner's Dilemma with payoff matrix $(R{=}3, S{=}0, T{=}5, P{=}1)$, against \emph{random} adversaries with defection probabilities of $0.5$ and $0.7$, with horizon \emph{disclosed} as 100 rounds; follow-up work in social dilemmas with multilingual variants reports similar prosocial bias under broadly comparable framings \cite{za2026towards}. The phenomenon is therefore reliably observed, but only in this restricted regime: known horizon, exploitable opponent, cooperative-payoff dominance. The picture inverts in genres outside this regime. In non-cooperative competitive settings such as iterated rock--paper--scissors, frontier LLMs play \emph{deeper} strategically than humans \cite{wang2026discovering}, and when game horizons are unknown, cooperation universalises across all languages \cite{buscemi2025fairgame, lore2024framing}, swamping any prosocial tilt. The prosocial default is therefore a property of cooperative-or-mixed-motive games with disclosed horizons against exploitable opponents --- not a uniform behavioural trait of LLMs. With that bound in place, the finding is still striking evolutionarily: thirty years of agent-based work following Axelrod \cite{axelrod1984evolution} suggest stable cooperation requires repeated interactions, reputation, and punishment. LLM cooperation is observed even when these scaffolds are weak, suggesting that cultural inheritance via pretraining produces a default that evolutionary game theory would not predict from interaction structure alone. We return to this in Section~\ref{sec:evolution}.

\subsection{Contradictory Evidence and Confounds}
\label{sec:confounds}

A responsible account of linguistic relativity in machines must foreground the evidence against it. Three findings are central. First, in a controlled cross-lingual study of persuasive co-writing for charitable appeals, \citeauthor{biswas2025persuasion} \cite{biswas2025persuasion} find that prior exposure to a Spanish LLM systematically reduces subsequent use of an English LLM, yet these usage patterns \emph{do not translate into aggregate differences in persuasiveness}; instead, recipients' \emph{beliefs} about whether the advertisement was AI-generated drive donation behaviour, particularly among Spanish-speaking female participants who downwardly adjusted donations when they suspected AI authorship. Second, when game horizons are unknown, cooperation \emph{universalises} across all languages \cite{buscemi2025fairgame, lore2024framing}: framing effects can fully dominate language effects. Third, apparent Vietnamese ``defection'' in low-stakes rounds sometimes reflects arithmetic miscomprehension under attenuated payoffs rather than strategic intent \cite{huynh2025fairgame}.

These findings bound the claim. Language effects exist, but they interact with framing, stakes, communication channels and task structure; they do not operate as a fixed cultural dial. \citeauthor{madmoun2025communication} \cite{madmoun2025communication} make this concrete: introducing a one-word communication channel raises Stag-Hunt cooperation from $0\%$ to $96.7\%$ across LLM agents irrespective of language, while a curriculum biased toward defection induces ``learned pessimism'' that lowers Public-Goods cooperation by $27.4\%$. Communication-channel presence and pretraining-curriculum balance can therefore swamp any language signal. Confounds are equally serious: tokeniser quality differs sharply across languages \cite{yong2024lowresource}, RLHF safety coverage is skewed \cite{ouyang2022training, yong2025multilingualsafety}, training corpora are unevenly distributed, and decoding hyperparameters interact with language-specific token statistics. Any causal claim about ``language X produces behaviour Y'' must be disentangled from these.

\section{Mechanism: Semantic Hub and Translation Barrier}
\label{sec:mechanism}

Two complementary accounts from mechanistic interpretability provide the machine-behaviour explanation for these patterns, and crucially for reviewers, dispense with the need for cultural-essentialist reasoning.

The \emph{Semantic Hub Hypothesis} \cite{wu2024semantichub}, supported by logit-lens decoding \cite{wendler2024llamasenglish}, holds that multilingual LLMs project heterogeneous linguistic inputs into a centralised, English-dominant representational hub in intermediate layers. Causal evidence comes from Activation Addition (ActAdd) experiments \cite{turner2023actadd}: \citeauthor{wu2024semantichub} \cite{wu2024semantichub} show that steering vectors derived in English (e.g., ``Good''$-$``Bad'' for sentiment) successfully shift outputs in Spanish and Chinese inputs, often as effectively as native-language steering, indicating that the shared representation space is not a vestigial byproduct but is causally engaged during cross-lingual processing. The corollary --- that non-English inputs are routed through English-anchored intermediate states --- predicts behavioural divergences when those intermediate states are sparsely populated for under-represented languages.

The complementary \emph{Translation Barrier Hypothesis} \cite{bafna2025translationbarrier} locates many multilingual failures not in upstream reasoning (which proceeds adequately in the hub) but in \emph{late-layer language realisation}: the model reaches the correct conclusion internally and then fails while projecting into the target language's token space. This predicts two diagnostic categories: \emph{hub failures} (intermediate representations do not encode the correct intent) and \emph{realisation failures} (intermediate representations are correct but surface output is not). Language-aware interpretability tools such as Indic-TunedLens offer methodology for the distinction \cite{kumar2025indictunedlens}.

Together these accounts predict the observed pattern: the bimodal Arabic IPD distribution and Vietnamese arithmetic errors reflect hub-level representational sparsity, while clembench format-adherence failures in German and Italian plausibly reflect late-layer realisation bottlenecks \cite{beyer2024clembench, schlangen2025clembench}. The same architecture produces both phenomena. A critical corollary is that \emph{confidence calibration} degrades in parallel: \citeauthor{ahuja2022calibration} \cite{ahuja2022calibration} report that Expected Calibration Error rises from approximately $4.6\%$ in English to $23.1\%$ on parallel non-English prompts on the MEGA benchmark across XGLM-7.5B, BLOOM-7.1B and PaLM-540B with $15$-bin equal-width binning, and \citeauthor{zhou2025beyond} \cite{zhou2025beyond} reproduce the asymmetry on Llama-3 and Qwen-2 families and show that probing a late-intermediate layer (rather than the final layer) recovers part of the gap. The result is over-confident, under-reliable agents precisely in the linguistic contexts where their substantive reasoning is weakest.

\subsection{Alternative Architectural Accounts}\label{sec:alternatives}

We do not claim that the hub--barrier dichotomy is the only architectural account on offer; the multilingual-representation literature supplies at least three live alternatives. \emph{Language-specific microcircuits}: \citeauthor{pires2019multilingualbert} \cite{pires2019multilingualbert} showed that multilingual BERT supports zero-shot transfer between typologically distant languages but exhibits ``systematic deficiencies affecting certain language pairs'', suggesting partially separated language-specific subnetworks alongside any shared hub. \emph{Subspace interference}: representation-learning analyses of XLM-R \cite{conneau2020xlmr} and mT5 \cite{xue2021mt5} show that scale alleviates but does not eliminate cross-lingual interference, with low-resource languages drifting in directions that high-resource languages constrain. \emph{Latent-language priors}: \citeauthor{wendler2024llamasenglish} \cite{wendler2024llamasenglish} report that Llama-family decoders pass through a latent intermediate that is statistically closer to English than to the input language --- a stronger claim than ``shared hub'' and the natural quantitative target for the depth-resolved probes proposed in P2 (Section~\ref{sec:predictions}). The architectural literature is therefore not exhausted by Semantic Hub vs.\ Translation Barrier; pre-registered probes that adjudicate among hub, microcircuit, interference and latent-language accounts on the same items are a precondition for sharper claims.

A natural test is provided by \emph{language-prioritised} models trained primarily on a non-English corpus: CT-LLM, with $800$B Chinese versus $300$B English tokens \cite{du2024ctllm}, and TigerLLM for Bangla \cite{raihan2025tigerllm}. If the Semantic Hub asymmetry is principally a function of training-data dominance rather than of architecture, these models should partially \emph{invert} the asymmetry --- better calibration and stronger ToM in their priority languages than in English. This is a falsifiable, immediately runnable test of the cultural-inheritance reading of multilingual LLM behaviour (Section~\ref{sec:evolution}).

\subsection{From Representation to Behaviour: Closing the Causal Loop}\label{sec:repbehavior}

A reasonable critic will point out that everything in this section so far is \emph{representational}. We have evidence that intermediate representations differ across languages and that they can be manipulated; we have separate evidence that behaviour differs across languages. The link between the two is, as it stands, inferential. Several recent results in interpretability supply both the methodology and the precedent for closing this loop directly. \citeauthor{chen2023spatialcausal} \cite{chen2023spatialcausal} show that spatial representations in DeBERTa and GPT-Neo have \emph{causal} effects on next-token prediction, not merely correlational structure. \citeauthor{zhang2026worldmodels} \cite{zhang2026worldmodels} demonstrate that latent world-model state variables in IRIS and DIAMOND are linearly decodable and causally engaged, again via targeted intervention. Layer-resolved analyses --- LogitLens4LLMs \cite{zhenyu2025logitlens}, the four-stage trajectory of arithmetic processing in Llama-3 \cite{yan2025arithmetic}, and the early-layer-syntax / late-layer-semantics dissociation in GPT-2 causal reasoning \cite{lee2024causalpaths} --- supply the concrete probe-and-intervene methodology that the multilingual case can now reuse. The natural next experiment, which we encode in our extended P2 (Section~\ref{sec:predictions}), is direct: derive an activation-steering vector from English IPD traces of cooperative versus defective play, apply it to the same model running an Arabic or Vietnamese IPD prompt, and measure whether the resulting strategy distribution shifts in the predicted direction. A positive result would convert the Semantic Hub from a representational curiosity into a behavioural lever; a negative result would force the field to look beyond the hub for the causal mechanism that the linguistic-relativity finding requires.

Two practical interventions are also worth noting because they bear on the governance discussion in Section~\ref{sec:society}: \emph{language adapters} that allocate per-language low-rank parameters on top of a shared backbone, and \emph{instruction-tuning variations} that explicitly cover under-represented languages, both demonstrably reduce calibration gaps and downstream behavioural divergence \cite{xue2021mt5, conneau2020xlmr}. These are not yet integrated into multi-agent evaluation pipelines; they should be.

\section{Theory of Mind Across Languages}
\label{sec:cognition}

Cooperation and deception both require \emph{theory of mind} (ToM): the capacity to model the mental states of other agents \cite{premack1978tom, deweerd2013bayesiantom}. In the cognitive-science canon, ToM is often the bottleneck behind cooperative achievement, and its cross-linguistic robustness matters for how LLM agents will perform in human multilingual groups.

A construct-validity caveat is in order before we proceed. The status of ToM benchmarks for LLMs is contested. \citeauthor{ullman2023tomfail} \cite{ullman2023tomfail} shows that small alterations to standard false-belief items collapse apparent ToM performance even in frontier models, suggesting the original successes may reflect surface pattern-matching rather than mental-state representation. \citeauthor{sap2022neuraltom} \cite{sap2022neuraltom} report that GPT-4 reaches only $\sim 60\%$ on SocialIQa and ToMi, well below human performance, and argue that scaling alone is insufficient. We therefore read cross-lingual ToM \emph{differences} as the more interpretable signal here: even if absolute ToM accuracy is contestable, the question of whether the same model exhibits asymmetric ToM \emph{across languages on parallel items} is well-defined and answerable independently of construct-validity disputes. Human evidence makes the question sharper. Developmental work suggests that bilingual children acquire false-belief understanding earlier than monolingual peers and pass standard ToM tasks at higher rates \cite{kovacs2009bilingual, goetz2003bilingual} --- a finding that links the executive demands of code-switching to the same control machinery that ToM recruits. The implicit prediction for machines is that ``multilingual'' systems should, if they genuinely manage multiple linguistic representations, show \emph{stronger} ToM than monolingual ones. The empirical pattern in LLMs is the opposite, and this contrast is informative: LLM multilinguality is implemented by sharing parameters across non-equivalently represented languages rather than by maintaining explicit cross-linguistic monitoring.

The XToM benchmark \cite{chan2025xtom} reveals what its authors describe as a ``pronounced dissonance'': while frontier models (GPT-4o, DeepSeek-R1) excel at multilingual language understanding, their ToM accuracy varies sharply across languages on the same belief-reasoning items. GPT-4o-11-20, for instance, exhibits an 11-percentage-point gap between Japanese and Chinese on XFANToM belief questions, and most evaluated LLMs --- with the notable exception of DeepSeek-R1 --- lack robust multilingual ToM reasoning despite achieving comparable accuracy on parallel fact-retrieval items. LLM-Hanabi \cite{chen2025llmhanabi} further establishes that \emph{first-order} ToM (interpreting another agent's intent) is the stronger predictor of cooperative success; when first-order ToM erodes in lower-resource languages, both cooperative \emph{and} deceptive reasoning suffer together. The result is counter-intuitive but coherent: a machine agent with weaker multilingual ToM is simultaneously a less effective collaborator and a less effective deceiver --- a combination that is safer on one axis but more brittle on another. We emphasise that this co-occurrence is currently \emph{correlational}: ToM-degraded behaviour and degraded commitment reliability are observed in the same languages, but the causal route --- whether ToM is the mediating bottleneck, a downstream symptom of the same hub sparsity, or an independent failure mode --- has not been adjudicated by mediation analysis or targeted ToM-impairment interventions. This is a clean empirical question that the field has the methodological tools to answer.

The intent--action gap provides a second cognitive measurement point. The MAST taxonomy, built from over 1,600 execution traces across seven multi-agent frameworks \cite{cemri2025mast}, categorises 14 failure modes into Specification/Design Issues ($41.77\%$), Inter-Agent Misalignment ($36.94\%$), and Task Verification ($21.30\%$). The core failures --- reasoning-action mismatches and unwarranted conversation resets --- are cognitive in nature, and M3-BENCH's \emph{commitment-reliability} metric shows that promise-keeping drops significantly in mixed-motive multilingual games \cite{xie2026m3bench}.

SPIN-Bench sharpens the picture by decoupling planning from social intelligence \cite{yao2025spinbench}: o-series reasoning models (o1, o1-preview, o1-mini) lead formal PDDL classical planning, yet \emph{introducing negotiation phases in Diplomacy degrades} the chain-of-thought coherence of the same strong planners, with o1's deal acceptance rate plateauing around 0.67 and other strong planners performing markedly worse on social metrics. Planning and ToM are distinct faculties that do not scale together. From a machine-behaviour standpoint, this mirrors the evolutionary-psychology distinction between ``Machiavellian intelligence'' and domain-general problem solving \cite{byrne1988machiavellian}, and suggests that cross-lingual transfer may be uneven across these faculties in ways that demand separate measurement.

\subsection{Moderating Factors Beyond Language}\label{sec:moderators}

Three recent studies establish that language is one modulator among several whose interactions deserve joint study. \citeauthor{ong2025personalitysteering} \cite{ong2025personalitysteering} show that representation-engineering steering of Big-Five Agreeableness and Conscientiousness reliably modulates Iterated Prisoner's Dilemma cooperation while increasing exploitation susceptibility --- direct evidence that internal representations \emph{can} causally steer cooperative play, which strengthens the more specific claim that \emph{language-linked} representations do so. \citeauthor{nishimoto2026commdelay} \cite{nishimoto2026commdelay} document a U-shaped relationship between communication delay and cooperation in continuous Prisoner's Dilemma settings, with agents spontaneously beginning to exploit slower partners; infrastructure can therefore reshape behaviour without any change in language. \citeauthor{jiang2025trustcollides} \cite{jiang2025trustcollides} show that human cooperation against an agent in repeated Prisoner's Dilemma depends on the \emph{declared identity} of that agent (human, rule-based AI, LLM); language is one of several signals that frame agent identity in human eyes. Any clean test of linguistic-relativity claims must hold these three moderators --- internal-trait steering, communication-channel infrastructure and partner-identity framing --- explicitly constant.

\section{Evolutionary Perspectives on Machine Cooperation}
\label{sec:evolution}

The ``nicer than human'' finding (Section~\ref{sec:relativity}) demands an evolutionary-game-theoretic explanation, and connecting LLM behaviour to this tradition is where the machine-behaviour research programme is least developed and most promising. Two intellectual lineages are immediately relevant. \citeauthor{han2022emergent} \cite{han2022emergent} has argued that evolutionary game theory provides the natural lens for studying emergent behaviour in multi-agent systems, including agents whose strategies arise from learning rather than mutation; \citeauthor{perc2017statistical} \cite{perc2017statistical} have shown that the statistical physics of cooperation --- spatial reciprocity, network reciprocity, phase transitions in cooperative regimes --- generalises beyond human populations to any well-defined population of strategic interactors. We propose that multilingual LLM populations satisfy the conditions of both frames.

\subsection{Axelrod's Legacy in a New Substrate}

Axelrod's computer tournaments established that cooperation can emerge among self-interested agents when interactions are iterated, payoffs favour mutual gains, and simple reciprocity (Tit-for-Tat) is available \cite{axelrod1984evolution, nowak2006five}. Subsequent work enriched this with reputation mechanisms, punishment, spatial structure, and indirect reciprocity \cite{nowak1998evolution, panchanathan2004indirect}. In all of these, cooperation is an \emph{achievement}, selected over a history of interactions.

LLM cooperation does not fit this mould. Models arrive at their first interaction \emph{pre-cooperative}: the pretraining and RLHF pipelines act as a form of \emph{cultural inheritance} that delivers a prosocial disposition before any game is played. This is most naturally read as a Boydian--Richersonian cultural-evolution process \cite{boyd1985culture, richerson2005notbygenes, henrich2016secret}: pretraining corpora are cumulative cultural products of millions of human authors, and RLHF injects an additional layer of selected human preferences. The resulting agents exhibit \emph{culturally inherited cooperation}, not evolved reciprocity.

This reframes the linguistic-relativity finding. If cooperation is culturally inherited via text, then languages with different textual corpora should deliver different cooperative priors --- exactly the pattern observed in Table~\ref{tab:strategy}. The prediction is testable: languages whose pretraining corpora are more heavily weighted toward cooperative, prosocial discourse (educational and informational genres, perhaps) should yield more cooperative play, while corpora dominated by adversarial or commercial text should yield more competitive defaults. Existing corpus-audit methods \cite{gao2020pile, luccioni2021what} provide the machinery; the hypothesis has not yet been tested directly.

A specific empirical concern follows. Pretraining corpora are massively skewed toward Western, Educated, Industrialised, Rich, Democratic (WEIRD) sources, yet \citeauthor{henrich2010weirdest} \cite{henrich2010weirdest} and \citeauthor{muthukrishna2020beyondweird} \cite{muthukrishna2020beyondweird} have shown that WEIRD populations are statistical \emph{outliers} on traits the cooperation literature treats as fundamental: individualism, analytic thinking, impersonal prosociality, sensitivity to fair-procedure framings. If LLM cooperative defaults are inherited from these corpora, they encode WEIRD-specific norms while presenting as universal. Multilingual deployment therefore risks not just performance gaps but normative imposition: a Swahili-language LLM agent defaulting to WEIRD prosociality may produce outcomes that are foreign to the cultural context it serves, while the same WEIRD prosociality is invisible (because culturally native) when deployed in English.

\subsection{Cultural Evolution of Machine Behaviour}

A further prediction follows. If LLM populations are linked --- via shared benchmarks, shared model families, and an increasingly common pipeline of distillation and self-training on model outputs --- then machine behaviour is itself subject to cultural evolution. Model-to-model transmission (distillation, synthetic-data training) is a form of \emph{high-fidelity oblique inheritance}; RLHF preferences are a form of \emph{selection}. The long-horizon question is whether machine behavioural traits drift in directions that human selection actively desires, or in directions driven by the structure of the transmission process itself --- a classical cultural-evolution question that the machine-behaviour programme is now equipped to ask empirically.

Language interacts with this process. If the transmission bottleneck is largely English --- because English is the lingua franca of benchmarks, distillation targets, and RLHF annotator pools --- then non-English behavioural traits may drift independently, neither as evolved nor as selected, simply as under-observed side effects. The Semantic Hub Hypothesis supplies the architectural substrate for this prediction \cite{wu2024semantichub}. Empirically, the substantial English--Arabic performance gap reported in TCG-Bench at the 32B-parameter scale \cite{wang2025tcgbench} is consistent with a drift that scale alone does not correct.

A clean test of the cultural-inheritance reading nonetheless requires disentangling three confounded sources of policy: the pretraining corpus, the RLHF policy prior, and decoding-time sampling. Existing studies typically vary one and confound the rest. The cleanest design is a $2 \times 2 \times 2$ ablation across (i) pretraining-language balance (matched vs.\ English-skewed); (ii) RLHF-annotator-language balance; (iii) decoding temperature/top-$p$, evaluated on a fixed multilingual IPD protocol. Until such a study exists, ``culturally inherited cooperation'' is the best-supported but not the only available reading of the evidence.

\subsection{Trust, Reciprocity, and Commitment}

Four concepts meet in this section. \emph{Trust} and \emph{reciprocity} are operationalised in commitment reliability \cite{xie2026m3bench}: when agents announce cooperation and then defect, trust among machine and human counterparts erodes. \emph{Cooperation} is the behavioural output measured in Table~\ref{tab:strategy} and the ``nicer than human'' finding \cite{fontana2024nicer}. \emph{Cultural evolution} \cite{boyd1985culture, henrich2016secret} is the transmission mechanism producing these behaviours without in-game selection, and \emph{theory of mind} is the cognitive prerequisite that limits how reliably trust and reciprocity can be implemented (Section~\ref{sec:cognition}). Framing LLM social behaviour in these terms makes it commensurable with a century of work in biology, anthropology, economics, and the statistical physics of cooperation \cite{perc2017statistical}, rather than treating machines as a brand-new object requiring a brand-new vocabulary.

\section{Societal Implications and Open Questions}
\label{sec:society}

Multilingual LLM agents are not a laboratory curiosity. By late 2025, Jacaranda Health's UlizaLlama provides health-information assistance in Swahili, Hausa, Yoruba, isiXhosa and isiZulu; Lelapa AI's InkubaLM serves several African low-resource languages; Cohere's Aya Expanse covers 23 languages including Vietnamese; and India's Bhashini and analogous national pipelines route public-sector queries through multilingual LLMs that millions of citizens depend on for health, finance, and legal information. Stanford's HAI Policy programme has explicitly mapped the ``language gap'' as a structural challenge for low-resource deployment \cite{stanfordhai2024languagegap}, and recent surveys document that even well-funded multilingual models remain markedly less safe and less calibrated outside English \cite{yong2025multilingualsafety, m-alert2025, wang2024evaluating}. If language modulates machine behaviour as Sections~\ref{sec:relativity}--\ref{sec:evolution} suggest, three societal consequences warrant attention.

\emph{Unequal quality of machine-mediated institutions.} If confidence calibration degrades fivefold outside English \cite{zhou2025beyond}, public-facing LLM agents in the Global South will systematically offer more over-confident, less reliable advice to populations with fewer alternatives. This is not a bug of individual models but a property of the pipeline.

\emph{Linguistically uneven attack surfaces.} Code-Switching Red-Teaming achieves a $46.7\%$ increase in attack success by mixing languages at the token level \cite{yoo2025csrt}; Agent-in-the-Middle attacks exploit inter-agent channels in multilingual multi-agent systems \cite{he2025agentmiddle}. Adversaries gain disproportionate leverage in exactly the linguistic contexts where safety data is sparsest. This is a safety externality of multilingual deployment, and the regulatory question --- how to incentivise safe deployment behaviour without producing perverse over-regulation --- is itself an evolutionary-game-theoretic problem \cite{han2022pledges}.

The hub--barrier framework helps unpack \emph{why} code-switching is so effective as an attack vector. Two predictions diverge here. \emph{Pivoting reading}: intra-utterance mixing routes the prompt through higher-resource (often English) intermediate states, so realisation succeeds even when the surface language is low-resource --- which would predict that code-switching \emph{reduces} certain failure modes, not just amplifies them. \emph{Attack-surface reading}: code-switching evades safety filters because filters are trained on monolingual surface forms while the hub processes the meaning regardless --- which predicts a pure increase in attack success, with no benign correlate. The empirical record so far supports the second reading more strongly than the first: CSRT raises attack success without a documented benign benefit. But the two predictions are not mutually exclusive, and a controlled study that pre-registers both an attack-success measure and a task-completion measure on the same code-switched prompts would adjudicate them. We flag this as a near-term experimental opportunity rather than a settled question.

\emph{Co-evolution with human multilingual communities.} LLM outputs re-enter human discourse, and thence future pretraining corpora. If the machines' language-modulated behavioural defaults shape the text humans produce --- in business email, in customer interactions, in education --- then the cultural-inheritance loop we described in Section~\ref{sec:evolution} becomes bidirectional. What this means for institutional norms, linguistic diversity, and trust in machine intermediaries is a research agenda, not a solved problem.

\subsection{Concrete Governance Levers}\label{sec:governance}

The prescriptive value of the framework above lies in three concrete governance levers, each immediately actionable.

\emph{Benchmark diversification.} Existing safety benchmarks (M-ALERT \cite{m-alert2025}, multilingual safety surveys \cite{yong2025multilingualsafety, wang2024evaluating}) should be paired with \emph{interactive} multi-agent benchmarks (clembench \cite{schlangen2025clembench}, MultiAgentBench \cite{nie2025multiagentbench}, TCG-Bench \cite{wang2025tcgbench}). The pairing matters: language-induced behavioural gaps and language-induced safety gaps should be reported jointly rather than as two separate literatures.

\emph{Multilingual RLHF coverage standards.} Annotator-pool composition by language and dialect should be reported on every model card. We propose, as a starting point for community discussion, a minimum coverage threshold of $\geq 5\%$ of RLHF preference labels per language for any model deployed across more than five official languages of operation, with documented dialectal stratification (urban vs.\ rural, standard vs.\ vernacular) for at least the highest-population dialect within each language. These numbers should be treated as a pragmatic floor, not a target; safety-critical deployments may require substantially more.

\emph{Audit standards for multi-agent deployments.} Code-switching adversarial probes \cite{yoo2025csrt} and Agent-in-the-Middle threat models \cite{he2025agentmiddle} should become a default test bed for any multilingual multi-agent product, in the same way that prompt-injection probes have become a default for single-agent deployments. None of these is a research breakthrough; what is missing is the institutional infrastructure that makes them routine.

\subsection{Four Testable Predictions}
\label{sec:predictions}

We close with four predictions formulated to be falsifiable with currently available methods. Each is a wager on the framework above.

\emph{P1 --- Corpus--cooperation correspondence.} Holding model architecture and RLHF pipeline fixed, the rank order of cooperation rates across languages in iterated social dilemmas should track the rank order of \emph{cooperative-discourse density} in those languages' pretraining corpora. We operationalise cooperative-discourse density as the per-million-token frequency of $n$-gram patterns annotated for prosocial speech-acts (commitment, acknowledgement, mutual benefit, reciprocity offer) in a stratified random sample of pretraining shards, calibrated against a hand-annotated gold set ($\geq 5{,}000$ items per language, double-annotated, Cohen's $\kappa \geq 0.6$). Frontier-model pretraining corpora are typically not public, which appears to bar direct measurement; we recommend three workarounds, in order of cleanliness. (i) Use \emph{open-data models} whose pretraining corpora are documented (e.g., the Pile \cite{gao2020pile}, RedPajama, the CT-LLM corpus \cite{du2024ctllm}, the TigerLLM Bangla corpus \cite{raihan2025tigerllm}); these supply ground truth for closed models that share architectural lineage. (ii) Use \emph{proxy corpora} matched on Common-Crawl shard distribution and language sampling weights documented in model cards. (iii) Use \emph{controlled pretraining runs} on small models (e.g., $\sim 1$B parameters) where the cooperative-discourse density of each language can be intentionally varied. The prediction is falsified if the Spearman rank correlation between cooperative-discourse density and IPD cooperation rate, controlling for RLHF annotator-language proportions and decoding temperature, is not statistically distinguishable from zero across $\geq 8$ languages.

\emph{P2 --- Hub--barrier dissociation, with a behavioural readout.} The original P2 stops at representational diagnosis; the present version adds the causal-behavioural step that closes the loop (Section~\ref{sec:repbehavior}). Pre-registered linear probes at layers $\lfloor L/3 \rfloor$ and $\lfloor 2L/3 \rfloor$ --- chosen to span the early-to-mid and mid-to-late blocks where prior work locates representational and realisation phenomena, respectively \cite{wendler2024llamasenglish, bafna2025translationbarrier, zhenyu2025logitlens} --- should classify multilingual failures into hub-level and realisation-level categories \cite{kumar2025indictunedlens}, with proportions differing systematically across languages. Per-language argmax-depth probes detect language-specific peak loci. Then, critically, the \emph{behavioural step}: derive an activation-steering vector \cite{turner2023actadd} from English IPD traces of cooperative versus defective play, apply it to the same model running an Arabic, Vietnamese or Bangla IPD prompt, and measure the resulting change in strategy distribution. If the steering vector shifts the strategy distribution in the direction predicted by the cooperative--defective contrast it was derived from, the Semantic Hub becomes a behavioural lever in addition to a representational claim. A null result --- failures distributed identically across layers, and steering producing no measurable strategy shift --- would force the field to look beyond the hub for the causal mechanism that the linguistic-relativity finding requires.

\emph{P3 --- Cultural-evolution drift under distillation.} A controlled chain of distillation experiments (model A trained on synthetic data from model B, evaluated in IPD across languages) should show that behavioural traits propagate with high fidelity in English but degrade in lower-resource languages, reproducing a measurable bottleneck. If degradation does not occur, the inheritance-via-distillation account is wrong.

\emph{P4 --- Evaluation reliability ceiling.} For multilingual LLM-as-Judge \cite{fu2025multilingualjudge}, panels of $\geq 3$ heterogeneous judges combined with self-translation and checklist rubrics should raise Fleiss' $\kappa$ above $0.5$ (``moderate'' agreement). If the recovery saturates below this, quantitative cross-lingual machine-behaviour science requires bilingual human adjudication, not just better aggregation.

These predictions are not exhaustive. They are chosen because each, if confirmed, would reshape an existing line of NLP or alignment research; and each, if falsified, would force a clean retreat from a specific claim of this perspective.

\section{Conclusions}

We have argued that language is an active modulator of behaviour in LLM agents, and that understanding this modulation requires integrating three traditions: \emph{behavioural characterisation} in the style of Rahwan's machine-behaviour programme \cite{rahwan2019machinebehaviour}; \emph{mechanistic interpretability}, which supplies the Semantic Hub and Translation Barrier accounts \cite{wu2024semantichub, bafna2025translationbarrier}; and \emph{evolutionary, cultural-evolution, and statistical-physics theories of cooperation} \cite{axelrod1984evolution, boyd1985culture, henrich2016secret, han2022emergent, perc2017statistical}, which provide the vocabulary of trust, reciprocity, inherited cooperation, and transmission dynamics. The emergent phenomena --- a prosocial default that exceeds typical human cooperation rates yet erodes language-dependently, asymmetric theory-of-mind degradation, and compounding safety vulnerabilities via code-switching --- are not isolated benchmark artefacts but coordinated signatures of a multilingual machine ethology that is only beginning to be measured.

We have been careful to foreground contradictory evidence: a cross-lingual persuasion null finding, framing effects that dominate language effects, and unresolved confounds. The Algorithmic Sapir--Whorf Switch remains a working hypothesis, not a proven law. The four testable predictions of Section~\ref{sec:predictions} are intended as the next move: each is falsifiable with current methods, and together they convert this perspective from a synthesis into a research agenda. We see this agenda as continuous with the long-running engagement of the life and physical sciences with behaviour, cooperation and evolution --- one to which machine populations now offer a new, uncomfortably consequential empirical substrate.

\end{multicols}

\enlargethispage{20pt}

\ethics{This perspective synthesises existing literature and involves no new human-subjects research. Cross-linguistic behavioural divergences reported here reflect model--language interactions under specific training distributions, not properties of languages or of human cultures; readers are urged to communicate findings with this framing to avoid cultural essentialism.}

\dataccess{No new data were generated. All figures and numerical claims reference the cited primary sources.}

\aucontribute{D.S.D.M.: conceptualisation, literature synthesis, writing --- original draft. V.N.: supervision, writing --- review and editing. [Co-author TBC].}

\competing{The authors declare no competing interests.}

\funding{To be added.}

\ack{The authors thank colleagues at the Faculty of Information Technology, University of Science, VNU-HCM for helpful discussion.}

\disclaimer{Views expressed are those of the authors and not of any affiliated institution.}

\printbibliography

\end{document}
